注意由紧合或光滑基变换，这些连接函子满足自然的相容性条件，因此该极限确实有意义。
把上述余极限中的 $\on{P}(\Sht_{\mu}^{\loc(m, n)})$ 换成 $\on{D}(\Sht_{\mu}^{\loc(m, n)})$，也就定义了 $\on{D}(\Sht_{\bar k}^\loc)$，它带有反常 $t$-结构。

\begin{rmk}
\label{R: sheaves on Shtloc}
由 \eqref{R:pre of Sht} 与注 \ref{R: grp theoretical Sht}，可用更群论化的语言解释上述定义。即写作
\[\frac{LG}{\Ad_\sigma K}=\underrightarrow\lim_{\mu}\underleftarrow\lim_n \frac{\Gr_\mu^{(n)}}{\Ad_\sigma K},\]
从而定义
\[\on{P}(\Sht^\loc_{\bar k})=\underrightarrow\lim_{\mu,m,n}\on{P}(\frac{\Gr^{(n)}_\mu}{\Ad_\sigma K_m}).\]
\end{rmk}

对每个支配余权 $\mu$ 以及使 $m-n$ 为 $\mu$-大的数对 $(m,n)$，有自然的拉回函子
\begin{equation*}
\label{E:functor P(Hk) to P(S)fin}
\Phi^{\loc(m,n)}:=\Res^{m,n}_{m,0}: \on{P}_{K\otimes \bar k}(\Gr_\mu)\cong \on{P}(\Hk_\mu^{\loc(m)})\to \on{P}(\Sht_\mu^{\loc(m,n)}),
\end{equation*}
由 \eqref{true-to-trun} 与紧合光滑基变换，它与上述连接态射可交换。
取余极限，就得到良定义的函子
\begin{equation}
\label{E:functor P(Hk) to P(S)}
\Phi^{\loc}: \mathrm{P}_{K\otimes\bar k}(\Gr\otimes\bar k) \to \mathrm{P}(\Sht_{\bar k}^\loc).
\end{equation}

\begin{rmk}
\label{R: more object}
注意 $\mathrm{P}(\Sht_{\bar k}^\loc)$ 中有许多对象并不来自 $\on{P}_{K\otimes \bar k}(\Gr\otimes\bar k)$ 在上述函子 $\Phi^{\loc}$ 下的像。例如如例 \ref{Ex:mStau mn} 所述，$\Sht^{\loc(n,n)}_0=\bfB G(\mO/\varpi^n)$。因此有限群 $G(\mO/\varpi^n)$ 的每个表示 $\rho$ 都在 $\Sht^{\loc(n,n)}_0$ 上定义局部系统，从而定义 $\mathrm{P}(\Sht_{\bar k}^\loc)$ 中的对象。只要 $n>0$，该对象就不在 $\Phi^\loc$ 的本质像中。更一般地，对每个 $\mu$ 与 $G(\mO)$ 的每个表示 $\rho$，都有 $\mathrm{P}(\Sht_{\bar k}^\loc)$ 中的对象 $\mF_{\mu,\rho}$，来自近循环函子构造 \cite{GL}。
\end{rmk}

现在我们定义范畴 $\rmP^{\Hk}(\Sht^\loc_{\bar k})$。它与 $\rmP(\Sht_{\bar k}^\loc)$ 有相同的对象，只需按下述方式说明态射空间及其复合。
称 $\rmP(\Sht_{\bar k}^\loc)$（等价地 $\rmP^\mathrm{Hk}(\Sht^\loc_{\bar k})$）中的对象 $\mF$ 为连通的，若它来自余极限 \eqref{E:definition of P(Sloc)} 中某个 $(\mu,m, n)$ 对应的 $\rmP(\Sht_\mu^{\loc(m,n)})$。此时记 $\mF_\mu^{(m,n)}\in \rmP(\Sht_\mu^{\loc(m,n)})$ 为代表 $\mF$ 的唯一（在同构唯一意义下）对象。由于 $\rmP(\Sht_{\bar k}^\loc)$ 的每个对象都是连通对象的直和，只需定义两个连通对象 $\calF_1, \calF_2$ 之间的态射空间，再按线性性推广到一般对象。

回忆对七元组 $(\mu_1,\mu_2,\nu,m_1,n_1,m_2,n_2)$，若 $(m_1,n_1,m_2,n_2)$ 是 $(\mu_1+\nu,\nu)$-可接受且 $(\mu_2+\nu,\nu)$-可接受的，则由命题 \ref{P:pushforward is local} 有对应 \eqref{deompf1}。给定满足类似可接受条件的另一个七元组 $(\mu'_1,\mu'_2,\nu',m'_1,n'_1,m'_2,n'_2)\geq (\mu_1,\mu_2,\nu,m_1,n_1,m_2,n_2)$，有下图
\begin{equation*}
\label{E: loc Hecke corr changing level}
\begin{CD}
\Sht_{\mu_1}^{\loc(m_1,n_1)}@<\res^{m'_1,n'_1}_{m_1,n_1}<< \Sht_{\mu_1}^{\loc(m'_1,n'_1)}@>i_{\mu_1,\mu'_1}>> \Sht_{\mu'_1}^{\loc(m'_1,n'_1)}\\
@AAA@AAA@AAA\\
\Sht_{\mu_1\mid\mu_2}^{\nu,\loc(m_1,n_1)}@<\res<<\Sht_{\mu_1\mid\mu_2}^{\nu,\loc(m'_1,n'_1)}@>i>>\Sht_{\mu'_1\mid\mu'_2}^{\nu',\loc(m'_1,n'_1)}\\
@VVV@VVV@VVV\\
\Sht_{\mu_2}^{\loc(m_2,n_2)}@<\res^{m'_2,n'_2}_{m_2,n_2}<<\Sht_{\mu_2}^{\loc(m'_2,n'_2)}@>i_{\mu_2,\mu'_2}>>\Sht_{\mu'_2}^{\loc(m'_2,n'_2)}
\end{CD}
\end{equation*}
其中向左的箭头是等维的、完全光滑的，且相对维数相同；向右的箭头是闭嵌入。注意由命题 \ref{P:pushforward is local}，左上方块是笛卡儿的，而右上方块中所有箭头都完全紧合。

现在，对 $\on{P}^\Hk(\Sht^\loc_{\bar k})$ 的两个连通对象 $\mF_1,\mF_2$，把 $\mathrm{Mor}_{\rmP^\Hk(\Sht^\loc_{\bar k})} (\calF_1, \calF_2)$ 定义为
\begin{equation*}
\label{E: hom space of PcorrSht}
\varinjlim_{(\mu_1,\mu_2,\nu,m_1,n_1,m_2,n_2)} \mathrm{Corr}_{\Sht_{\mu_1|\mu_2}^{\nu, \loc(m_1,n_1)}} \big( (\Sht_{\mu_1}^{\loc(m_1,n_1)}, \calF_{1,\mu_1}^{(m_1,n_1)}),\,( \Sht_{\mu_2}^{\loc(m_2,n_2)},\calF_{2,\mu_2}^{(m_2,n_2)}) \big),
\end{equation*}
其中余极限取遍所有（带偏序的）满足上述可接受条件、且使 $\mF_{i,\mu_i}^{(m_i,n_i)}$ 代表 $\mF_i$ 的七元组，余极限的连接映射由 $i_!\circ\res^![((m'-n')-(m-n))\dim G]$ 给出，其中 $i_!$ 表示按 \eqref{ASS:pushforward correspondence} 沿 $i$ 的上同调对应前推，$\res^!$ 表示按 \eqref{ASS:smooth pullback correspondence2} 的上同调对应拉回。如 \cite[\S 5.4.1]{XZ} 所述，这些连接映射可以复合，该余极限良定义。

接着说明态射的复合。由线性性，只需对连通对象定义，并可设所给态射 $c_1: \calF_1 \to \calF_2$ 与 $c_2: \calF_2 \to \calF_3$ 都实现为对应（而非对应的线性组合）
{\footnotesize
\begin{align*}
c^{(\mu_1,\mu_2,\nu_1,m_1,n_1,m_2,n_2)}_1 \in \mathrm{Corr}_{\Sht_{\mu_1|\mu_2}^{\nu_1,\loc(m_1,n_1)}}\big( (\Sht_{\mu_1}^{\loc(m_1, n_1)}, \calF_{1,\mu_1}^{(m_1,n_1)}),\ (\Sht_{\mu_2}^{\loc(m_2, n_2)}, \calF_{2,\mu_2}^{(m_2,n_2)}) \big)&, \quad\textrm{以及}\\
c^{(\mu_2,\mu_3,\nu_2,m_2,n_2,m_3,n_3)}_2 \in \mathrm{Corr}_{\Sht_{\mu_2|\mu_3}^{\nu_2,\loc(m_2,n_2)}}\big( (\Sht_{\mu_2}^{\loc(m_2, n_2)},\ \calF_{2,\mu_2}^{(m_2,n_2)}) \,\ (\Sht_{\mu_3}^{\loc(m_3, n_3)}, \calF_{3,\mu_3}^{(m_3,n_3)})\big)&,
\end{align*}}
其中各七元组满足上述条件。
按上同调对应复合的形式体系（见定义 \ref{AD:correspondences}），复合 $c^{(\mu_2,\mu_3,\nu_2,m_2,n_2,m_3,n_3)}_2 \circ
c^{(\mu_1,\mu_2,\nu_1,m_1,n_1,m_2,n_2)}_1$ 支撑在如下对应上
\[
\Sht_{\mu_1}^{\loc(m_1,n_1)} \leftarrow \Sht_{\mu_1|\mu_2}^{\nu_1,\loc(m_1, n_1)} \times_{ \Sht_{\mu_2}^{\loc(m_2,n_2)}}
\Sht_{\mu_2|\mu_3}^{\nu_2,\loc(m_2, n_2)}  \to  \Sht_{\mu_3}^{\loc(m_3,n_3)}.
\]

然后把复合 $c_2\circ c_1$ 定义为从 $\calF_1$ 到 $\calF_3$ 的上同调对应，使 $(c_2\circ c_1)^{(\mu_1,\mu_3,\nu_1+\nu_2,m_1,n_1,m_3,n_3)}$ 由 $c^{(\mu_2,\mu_3,\nu_2,m_2,n_2,m_3,n_3)}_2 \circ
c^{(\mu_1,\mu_2,\nu_1,m_1,n_1,m_2,n_2)}_1$ 沿完全紧合态射
\[
\Sht_{\mu_1|\mu_2}^{\nu_1,\loc(m_1, n_1)} \times_{ \Sht_{\mu_2}^{\loc(m_2,n_2)}}
\Sht_{\mu_2|\mu_3}^{\nu_2,\loc(m_2, n_2)} \xrightarrow{\mathrm{Comp}^{\loc(m_1,n_1)}}\Sht_{\mu_1|\mu_3}^{\nu_1+\nu_2,\loc(m_1, n_1)}
\]
的前推（按 \eqref{ASS:pushforward correspondence} 的意义）给出。利用命题 \ref{P:pushforward is local} 与 \S \ref{Sec:cohomological correspondence} 中的若干引理，可以证明 $c_2\circ c_1$ 良定义，并满足自然的结合律 $(c_3\circ c_2) \circ c_1 \cong c_3 \circ (c_2 \circ c_1)$。至此 $\on{P}^{\Hk}(\Sht^\loc_{\bar k})$ 的构造完成。

我们描述 $\on{P}^{\on{Hk}}(\Sht^\loc_{\bar k})$ 中某些对象的自同态环。
对每个 $n>0$，记 $\on{reg}_n$ 为 $G(\mO/\varpi^n)$ 的正则表示。它在 $\on{Sht}^{\loc(n,n)}_{0}=\bfB G(\mO/\varpi^n)$ 上定义局部系统
$\delta_{\on{reg}_n}^{(n,n)}$（见注 \ref{R: more object}）。记 $\delta_{\on{reg}_n}$ 为 $\on{P}^{\Hk}(\Sht^\loc_{\bar k})$ 中相应的对象。
注意 $\delta_{\on{reg}_0}$ 含于函子 \eqref{E:functor P(Hk) to P(S)} 的像中，我们也把它记作 $\widetilde{\delta_{e}}$。
下面的结论是 \eqref{E:endo for unit} 与 \eqref{E:endo for reg} 的仿射类似物。
\begin{prop}
\label{P: endo of unit}
有典范同构
$$\End_{\on{P}^{\Hk}(\Sht^\loc_{\bar k})}(\delta_{\on{reg}_n})\cong C_c(K(n)\backslash G(F)/K(n),E)^{\on{op}},$$ 其中 $K(n)=K^{(n)}(k)$ 是第 $n$ 级主同余子群，$C_c(K(n)\backslash G(F)/K(n),E)$ 表示相应的 Hecke 代数。
特别地，$\End_{\on{P}^{\Hk}(\Sht^\loc_{\bar k})}(\widetilde{\delta_{e}})$ 同构于 $E$ 值球面 Hecke 代数 $H_G\otimes E=C_c(G(\mO)\backslash G(F)/G(\mO),E)$。
\end{prop}

\begin{rmk}
完全相同的定义给出 $\on{D}^\Hk(\Sht^{\loc}_{\bar k})$。此外，若 $\on{Iw}\subset K$ 是 Iwahori 子群，则可类似地定义带 Iwahori 层级结构的局部 Shtuka 的模空间，记作 $\Sht^{\loc}_{\on{Iw}}$。同样可以定义 Hecke 对应 $\Hk(\Sht^{\loc}_{\on{Iw}})$ 与范畴 $\on{D}^\Hk(\Sht^{\loc}_{\on{Iw},\bar k})$。如注 \ref{R:springer} 所述，若 $\cha E=0$，则它在幂等完备化之后等价于 $\on{D}^\Hk(\Sht^{\loc}_{\bar k})$。
\end{rmk}
